This chapter surveys the body of work that underpins the present research. It is organized into five areas: conversational product recommendation systems and their evolution, Retrieval-Augmented Generation (RAG) within the context of e-commerce, the application of agentic workflows for dynamic decision support, methods for resolving ambiguity in natural language queries, and the optimization of system workflows through tool-augmented execution versus direct model autonomy. The chapter concludes with a summary that identifies the specific gap this thesis addresses, particularly in the domain of personalized fashion suggestions for users with ill-defined or unclear intents.
\section{Conversational product recommendation systems and their evolution}

Traditional recommendation systems (RS) primarily utilize collaborative filtering and content-based approaches. While effective in high-data environments, these systems often lack transparency and do not allow for real-time user feedback during the discovery process. The development of Conversational Recommendation Systems (CRS) addressed these limitations by facilitating multi-turn dialogues. Early iterations relied on predefined decision trees and intent classification, which often resulted in rigid user experiences. Recent advancements in Large Language Models (LLMs) have enabled more sophisticated natural language interactions, though these models frequently require additional grounding to remain consistent with specific product inventories.

\section{Retrieval-Augmented Generation (RAG) within the context of e-commerce}

The integration of RAG has significantly mitigated the risk of generating inaccurate or non-existent product information—a common challenge when using generative models in retail. By linking generative architectures with external vector databases, RAG ensures that outputs are derived from verified catalog data. In e-commerce, this allows for semantic retrieval based on descriptive user needs rather than simple keyword matching. Current research emphasizes the necessity of grounding for system reliability; however, standard RAG implementations often fail when initial user queries are insufficient to trigger a precise retrieval.

\section{The application of Agentic workflows for dynamic decision support}

Recent shifts in AI architecture have introduced Agentic workflows, which incorporate autonomous reasoning and iterative planning into the system. In a dynamic decision support context, an Agent functions as a controller that evaluates the sufficiency of the provided information before proceeding. Frameworks such as ReAct (Reason + Act) enable systems to decompose complex user requests into manageable tasks. For fashion assistants, this enables the system to identify missing parameters in a query and strategically initiate follow-up interactions to refine the search space.

\section{Methods for Resolving Ambiguity in Natural Language Queries}

A persistent challenge in fashion retrieval is the presence of indeterminate user intent—cases where queries are abstract or linguistically vague. Research in Natural Language Understanding (NLU) has explored various strategies for disambiguation, including query expansion and context-aware rewriting. Modern approaches utilize LLMs to translate "lifestyle-based" descriptions into structured technical attributes. Despite these gains, maintaining a balance between autonomous inference and explicit user clarification remains a critical area of investigation for improving retrieval precision.

\section{The optimization of system workflows through tool-augmented execution versus direct model autonomy}

Current academic and industry discussions contrast the efficacy of direct model autonomy against tool-augmented execution. Direct autonomy, where the LLM is responsible for both reasoning and data output, is susceptible to stochastic errors and logical inconsistencies. Conversely, tool-augmented execution utilizes the LLM as an orchestrator for deterministic software components, such as database APIs, filtering scripts, and inventory management systems. This research posits that by delegating execution to specific tools, the system achieves higher reliability and operational accuracy, ensuring that recommendations are consistent with real-world availability.

\section{Summary and Research Gap}

While the literature extensively covers the components of RAG and Conversational AI, there is a distinct lack of research focused on Agentic systems optimized for the high-variance domain of fashion retrieval. Existing solutions often rely either on static filtering or unconstrained generative models. This thesis addresses this gap by developing Clothie. This production-oriented Agentic RAG system employs tool-augmented workflows to resolve ambiguous user queries while prioritising the \textbf{fastest practical response path} for user experience, thereby bridging the gap between indeterminate intent and precise, verifiable product suggestions.